The framework serves a defender: a detector developer or an analyst who maintains static detectors for executable
files and wants to know what they rely on. We state what the defender needs, which adversary the applications
anticipate, and what is out of scope.

\paragraph{The defender.} The defender owns the detectors and their feature extractors and can run them on any file.
To compute the provenance matrix $P$ exactly, the defender instruments the extractor's source code, once; without the
source, $P$ can be estimated from the extractor's outputs alone by occluding windows of bytes, at a cost in accuracy
that we measure (Section~\ref{sec:XML:CharacterizingP}). The transfers take an explanation as input and are only as
faithful as it is. The explainers we use need the access a model owner has: TreeSHAP reads the trees, and Integrated
Gradients reads the raw-byte network's gradients; KernelSHAP and occlusion need only queries. The defender also has
a labeled corpus of files on which to explain and audit the detectors.

\paragraph{The tasks.} We evaluate two security tasks, one per transfer. \textit{Disagreement triage}
(Section~\ref{sec:XML:Comparison}) is local: when two of the defender's detectors disagree on a file, decide whether
the flagging detector's evidence is content the other does not read or reads the other way, and find the bytes both
rely on. \textit{Reliance audit} (Section~\ref{sec:XML:Audit}) is global: find content that a raw-byte detector's
detections rest on across a corpus, so that the defender learns where the detector is cheap to evade before an
attacker does.

\paragraph{The adversary the audit anticipates.} A malware author wants a file that the detector flags to be
classified benign, and can change any bytes of the file, provided it still loads and runs. The author does not need
the detector's explanations or parameters; finding the content a detector relies on is exactly what the defender's
audit does first. We measure reliance by scrambling the bytes $P$ names for a feature, with an operator that keeps the
file parseable and its layout intact but does not verify that the program still runs. Our evasion rates therefore
measure how much of the detection rests on content an attacker could target, not a demonstrated attack: some of these
fields cannot be changed freely without breaking the program, and an attacker who chooses values rather than random
bytes could do better.

\paragraph{Out of scope.} Adversaries who attack the explanations themselves (e.g., by manipulating a model so that its
explanations look benign), the feature extractor, or the defender's training data. Our planted triggers
(Section~\ref{sec:XML:Setup}) are poisoning by construction, but we use them only as a known answer for evaluating
explanations, not as a threat we defend against. Dynamic analysis features, computed from a file's behavior rather
than its bytes, are also out of scope: their provenance would be a trace of execution, not a set of bytes.
